% =================================================================
\section{Calculus I: differentiation and its applications (Area of Study 3)}
\label{sec:diff}
% =================================================================

\begin{sdbox}
\textbf{Calculus covers} \emph{``graphical treatment of limits, continuity and differentiability of functions of a single real variable, and differentiation, anti-differentiation and integration of these functions''}. The differentiation dot points are:
\begin{itemize}
  \item deducing the graph of the derivative function from the graph of a given function
  \item derivatives of $x^n$ for $n\in\Q$, $e^x$, $\loge(x)$, $\sin(x)$, $\cos(x)$ and $\tan(x)$
  \item derivatives of $f(x)\pm g(x)$, $f(x)\times g(x)$, $\dfrac{f(x)}{g(x)}$ and $(f\circ g)(x)$, where $f$ and $g$ are polynomial, exponential, circular, logarithmic or power functions and transformations or simple combinations of these
  \item application of differentiation to graph sketching and identification of key features, including stationary points and points of inflection, and intervals over which a function is strictly increasing or strictly decreasing
  \item identification of local maximum/minimum values over an interval and application to solving optimisation problems in context, including identification of interval endpoint maximum and minimum values
\end{itemize}
\textbf{Key knowledge/skills:} the link between the graph of a function and its gradient function; the tangent at a point; how the sign and magnitude of $f'$ describe $f$; the sum, difference, chain, product and quotient rules (\byhand\ for simple combinations); derivatives of $f(ax+b)$ for $f=x^n,\sin,\cos,\tan,e^x,\loge x$ \byhand; \textbf{Newton's method}.
\end{sdbox}

\subsection{Limits, continuity and differentiability}
\label{sec:cont}

\begin{defn}{Definitions (graphical treatment)}
\begin{itemize}
  \item $\displaystyle\lim_{x\to a}f(x)=L$: $f(x)$ gets arbitrarily close to $L$ as $x$ approaches $a$ from both sides. The left limit $\lim_{x\to a^-}$ and right limit $\lim_{x\to a^+}$ must be equal. Limits at infinity describe asymptotes, e.g.\ $\lim_{x\to\infty}e^{-x}=0$.
  \item $f$ is \textbf{continuous at $x=a$} if $f(a)$ is defined, $\lim_{x\to a}f(x)$ exists, and they are equal. Informally, you can draw through $x=a$ without lifting your pen.
  \item $f$ is \textbf{differentiable at $x=a$} if $f'(a)=\lim_{h\to0}\dfrac{f(a+h)-f(a)}{h}$ exists: the graph is \emph{smooth} there (left and right gradients equal and finite).
  \item Differentiable $\Rightarrow$ continuous, but \textbf{continuous $\not\Rightarrow$ differentiable}.
  \item \textbf{Not differentiable} at: corners (e.g.\ at hybrid joins), cusps ($x^{2/3}$ at $0$), vertical tangents ($x^{1/3}$ at $0$), any discontinuity, and the \textbf{endpoints} of a closed domain. The domain of $f'$ excludes all of these, so $f'$ is defined on open intervals.
\end{itemize}
\end{defn}

\begin{center}
\begin{tikzpicture}
\begin{axis}[smallgraph, width=4.3cm, height=3.6cm, title={\scriptsize corner}, xmin=-1.3, xmax=3.3, ymin=-0.5, ymax=3.2, xtick=\empty, ytick=\empty]
  \addplot[plotA, domain=-1.2:1] {x^2+1};
  \addplot[plotA, domain=1:3.1] {3-x};
  \node[pt] at (axis cs:1,2) {};
\end{axis}
\end{tikzpicture}\hfill
\begin{tikzpicture}
\begin{axis}[smallgraph, width=4.3cm, height=3.6cm, title={\scriptsize cusp: $y=x^{2/3}$}, xmin=-2.2, xmax=2.2, ymin=-0.4, ymax=1.9, xtick=\empty, ytick=\empty]
  \addplot[plotA, domain=-2:2] {abs(x)^(2/3)};
\end{axis}
\end{tikzpicture}\hfill
\begin{tikzpicture}
\begin{axis}[smallgraph, width=4.3cm, height=3.6cm, title={\scriptsize vertical tangent: $y=x^{1/3}$}, xmin=-2.2, xmax=2.2, ymin=-1.5, ymax=1.5, xtick=\empty, ytick=\empty]
  \addplot[plotA, domain=-2:2] {sign(x)*abs(x)^(1/3)};
\end{axis}
\end{tikzpicture}\hfill
\begin{tikzpicture}
\begin{axis}[smallgraph, width=4.3cm, height=3.6cm, title={\scriptsize jump discontinuity}, xmin=-1.3, xmax=2.6, ymin=-1.2, ymax=3.5, xtick=\empty, ytick=\empty]
  \addplot[plotA, domain=-1.2:1] {x};
  \addplot[plotA, domain=1:2.4] {x+1};
  \node[opt] at (axis cs:1,1) {};
  \node[pt] at (axis cs:1,2) {};
\end{axis}
\end{tikzpicture}
\end{center}

\begin{example}{making a hybrid function smooth}
Find $a$ and $b$ so that $f(x)=\begin{cases}ax^2+b, & x\le1\\ 4\sqrt{x}, & x>1\end{cases}$ is differentiable at $x=1$.\par
\textbf{Continuity:} $a(1)^2+b=4\sqrt1\Rightarrow a+b=4$.\quad
\textbf{Equal gradients:} left $2ax\big|_{x=1}=2a$; right $\ddx\left(4x^{1/2}\right)=2x^{-1/2}\big|_{x=1}=2$. So $2a=2$, giving $a=1$ and $b=3$.
\end{example}

\begin{keyres}{Rates of change and first principles (assumed from Units 1 \& 2)}
\begin{itemize}
  \item \textbf{Average rate of change} of $f$ over $[a,b]$: $\dfrac{f(b)-f(a)}{b-a}$, the gradient of the chord (secant).
  \item \textbf{Instantaneous rate of change} at $x=a$: $f'(a)$, the gradient of the tangent.
  \item First principles: $f'(x)=\displaystyle\lim_{h\to0}\frac{f(x+h)-f(x)}{h}$. E.g.\ for $f(x)=x^2$: $\dfrac{(x+h)^2-x^2}{h}=2x+h\to2x$.
\end{itemize}
\end{keyres}

% -----------------------------------------------------------------
\subsection{Standard derivatives and the rules}
\label{sec:rules}

\begin{cols}[0.47]
\begin{keyres}{Standard derivatives}
\renewcommand{\arraystretch}{1.45}
\begin{tabular}{@{}l@{\quad}l@{}}
$f(x)$ & $f'(x)$ \\ \hline
$x^n$, $n\in\Q$ & $nx^{n-1}$ \\
$(ax+b)^n$ & $an(ax+b)^{n-1}$ \\
$e^{x}$ & $e^{x}$ \\
$e^{ax+b}$ & $ae^{ax+b}$ \\
$\loge(x)$, $x>0$ & $\dfrac1x$ \\
$\loge(ax+b)$ & $\dfrac{a}{ax+b}$ \\
$\sin(ax+b)$ & $a\cos(ax+b)$ \\
$\cos(ax+b)$ & $-a\sin(ax+b)$ \\
$\tan(ax+b)$ & $\dfrac{a}{\cos^2(ax+b)}=a\sec^2(ax+b)$ \\
$a^x=e^{x\loge a}$ & $(\loge a)\,a^x$ \\
\end{tabular}
\par\smallskip
{\small Note: $\ddx\loge(kx)=\dfrac1x$, since $\loge(kx)=\loge k+\loge x$.}
\end{keyres}
\nextcol[0.505]
\begin{keyres}{The rules (with $u=u(x)$, $v=v(x)$)}
\begin{itemize}
  \item Sum/difference: $(f\pm g)'=f'\pm g'$; \ constant multiple $(kf)'=kf'$.
  \item \textbf{Product rule:} $\dfrac{d}{dx}(uv)=u\dfrac{dv}{dx}+v\dfrac{du}{dx}$
  \item \textbf{Quotient rule:} $\dfrac{d}{dx}\!\left(\dfrac uv\right)=\dfrac{v\frac{du}{dx}-u\frac{dv}{dx}}{v^2}$
  \item \textbf{Chain rule:} $\dfrac{dy}{dx}=\dfrac{dy}{du}\cdot\dfrac{du}{dx}$, i.e.\ $\big(f(g(x))\big)'=f'(g(x))\,g'(x)$
\end{itemize}
\textbf{Chain rule patterns} (learn these):
\begin{itemize}
  \item $\ddx[f(x)]^n=n\,f'(x)\,[f(x)]^{n-1}$
  \item $\ddx e^{f(x)}=f'(x)\,e^{f(x)}$
  \item $\ddx\loge(f(x))=\dfrac{f'(x)}{f(x)}$
  \item $\ddx\sin(f(x))=f'(x)\cos(f(x))$
  \item $\ddx\cos(f(x))=-f'(x)\sin(f(x))$
  \item $\ddx\tan(f(x))=\dfrac{f'(x)}{\cos^2(f(x))}$
\end{itemize}
\end{keyres}
\end{cols}

\begin{example}{differentiation \byhand\ (typical Exam 1)}
\renewcommand{\arraystretch}{1.95}
\begin{tabularx}{\linewidth}{@{}>{$}l<{$}L@{}}
y=(2x-1)^5 & $\dydx=5(2x-1)^4\times2=10(2x-1)^4$ \hfill (chain) \\
y=\sqrt{3x^2+1}=(3x^2+1)^{1/2} & $\dydx=\frac12(3x^2+1)^{-1/2}\times6x=\dfrac{3x}{\sqrt{3x^2+1}}$ \hfill (chain) \\
y=x^2e^{3x} & $\dydx=2x\,e^{3x}+x^2\cdot3e^{3x}=xe^{3x}(3x+2)$ \hfill (product) \\
y=\dfrac{\loge x}{x} & $\dydx=\dfrac{x\cdot\frac1x-\loge x\cdot1}{x^2}=\dfrac{1-\loge x}{x^2}$ \hfill (quotient) \\
y=\dfrac{x}{x^2+1} & $\dydx=\dfrac{(x^2+1)-x(2x)}{(x^2+1)^2}=\dfrac{1-x^2}{(x^2+1)^2}$ \hfill (quotient) \\
y=\cos(x^2) & $\dydx=-2x\sin(x^2)$ \hfill (chain) \\
y=\loge(\cos x) & $\dydx=\dfrac{-\sin x}{\cos x}=-\tan x$ \hfill (chain, log pattern) \\
y=e^{\sin x} & $\dydx=\cos x\,e^{\sin x}$ \hfill (chain) \\
y=\tan(2x) & $\dydx=\dfrac{2}{\cos^2(2x)}$ \\
f(x)=x\sin(2x) & $f'(x)=\sin(2x)+2x\cos(2x)$, so $f'\!\left(\frac{\pi}{4}\right)=\sin\frac{\pi}{2}+\frac{\pi}{2}\cos\frac{\pi}{2}=1$ \\
\end{tabularx}
\end{example}

\begin{tip}[Make differentiation easier]
\begin{itemize}
  \item Rewrite roots and fractions as powers first: $\dfrac{3}{\sqrt x}=3x^{-1/2}$, \ $\dfrac{2}{(1-x)^3}=2(1-x)^{-3}$.
  \item Simplify logs first: $\loge\!\left(\dfrac{x^2}{\sqrt{x+1}}\right)=2\loge|x|-\frac12\loge(x+1)$, so the derivative is $\dfrac2x-\dfrac{1}{2(x+1)}$. (Write $\loge|x|$, or state $x>0$: $\loge(x^2)$ is defined for $x<0$ but $2\loge x$ is not.)
  \item Factorise the answer (e.g.\ $xe^{3x}(3x+2)$). You will usually need $f'(x)=0$ next, and the factor form makes that easy.
  \item ``Show that $f'(x)=\dots$'': show the rule you used and at least one intermediate line.
\end{itemize}
\end{tip}

% -----------------------------------------------------------------
\subsection{Tangents and normals}
\label{sec:tangent}

\begin{cols}
\begin{keyres}{Equations}
At the point $(a,f(a))$:
\begin{itemize}
  \item \textbf{Tangent:} $y-f(a)=f'(a)(x-a)$.
  \item \textbf{Normal} (perpendicular to the tangent): $y-f(a)=-\dfrac{1}{f'(a)}(x-a)$ if $f'(a)\ne0$; if $f'(a)=0$ the normal is the vertical line $x=a$.
  \item Angle of the tangent with the positive $x$-axis: $\tan\theta=f'(a)$.
  \item ``Parallel to $y=mx+c$'': solve $f'(x)=m$.
  \item Linear approximation: $f(a+h)\approx f(a)+hf'(a)$ (the tangent line near $x=a$).
\end{itemize}
\end{keyres}
\nextcol
\begin{example}{tangent, normal, external point \byhand}
\textbf{(a)} $y=x^3-2x$ at $x=1$: $y=-1$, $\dydx=3x^2-2=1$.\\ Tangent $y=x-2$; normal $y=-x$.\par
\textbf{(b)} Tangents to $y=x^2$ that pass through $(1,-3)$. The tangent at $(a,a^2)$ is $y=2ax-a^2$. Substitute $(1,-3)$: $-3=2a-a^2\Rightarrow a^2-2a-3=0\Rightarrow a=3$ or $a=-1$. The tangents are $y=6x-9$ and $y=-2x-1$.
\end{example}
\end{cols}

% -----------------------------------------------------------------
\subsection{The graph of $f'$ from the graph of $f$}
\label{sec:gradgraph}

\begin{cols}[0.53]
\begin{keyres}{Reading the gradient}
\small
\begin{tabular}{@{}p{3.75cm}p{4.45cm}@{}}
\textbf{On the graph of $f$} & \textbf{On the graph of $f'$} \\ \hline
stationary point & $x$-intercept \\
$f$ increasing & above the $x$-axis ($f'>0$) \\
$f$ decreasing & below the $x$-axis ($f'<0$) \\
point of inflection & turning point of $f'$ \\
stationary inflection & touches the $x$-axis \\
steeper & larger $|f'(x)|$ \\
vertical asymptote $x=a$ & vertical asymptote $x=a$ (usually) \\
horizontal asymptote $y=c$ & $f'\to0$ (asymptote $y=0$) \\
polynomial of degree $n$ & polynomial of degree $n-1$ \\
corner, cusp, endpoint & $f'$ undefined there (open circles) \\
\end{tabular}
\end{keyres}
\begin{tip}[Sketching $f'$ from $f$]
Mark the $x$-values of the stationary points of $f$ (these are zeros of $f'$) and of any points of inflection (turning points of $f'$). Decide the sign of $f'$ on each interval, then join up smoothly. Exam questions often give the graph of $f'$ and ask about $f$: then read the same table from right to left.
\end{tip}
\nextcol[0.445]
\centering
\begin{tikzpicture}
\begin{axis}[mygraph, width=6.4cm, height=4.4cm, xmin=-2.3, xmax=2.5, ymin=-3.2, ymax=3.2, xtick={-1,1}, ytick=\empty, cornertitle, title={$f(x)=x^3-3x$}]
  \addplot[plotA, domain=-2.1:2.1] {x^3-3*x};
  \node[pt] at (axis cs:-1,2) {}; \node[lab, above] at (axis cs:-1,2) {max};
  \node[pt] at (axis cs:1,-2) {}; \node[lab, below] at (axis cs:1,-2) {min};
  \node[pt] at (axis cs:0,0) {}; \node[lab, right] at (axis cs:0.1,0.35) {inflection};
  \draw[dashed, black!40] (axis cs:-1,-3.2) -- (axis cs:-1,3.2);
  \draw[dashed, black!40] (axis cs:1,-3.2) -- (axis cs:1,3.2);
\end{axis}
\end{tikzpicture}

\begin{tikzpicture}
\begin{axis}[mygraph, width=6.4cm, height=4.4cm, xmin=-2.3, xmax=2.5, ymin=-3.8, ymax=6.5, xtick={-1,1}, ytick=\empty, cornertitle, title={$f'(x)=3x^2-3$}]
  \addplot[plotB, domain=-2.1:2.1] {3*x^2-3};
  \node[pt] at (axis cs:-1,0) {}; \node[pt] at (axis cs:1,0) {};
  \node[pt] at (axis cs:0,-3) {}; \node[lab, right] at (axis cs:0,-3) {$(0,-3)$ min of $f'$};
  \draw[dashed, black!40] (axis cs:-1,-3.8) -- (axis cs:-1,6.5);
  \draw[dashed, black!40] (axis cs:1,-3.8) -- (axis cs:1,6.5);
  \node[lab] at (axis cs:-1.75,1.4) {$+$};
  \node[lab] at (axis cs:1.75,1.4) {$+$};
  \node[lab] at (axis cs:0.5,-1.2) {$-$};
\end{axis}
\end{tikzpicture}
\end{cols}

% -----------------------------------------------------------------
\subsection{Stationary points, points of inflection, increasing and decreasing}
\label{sec:sp}

\begin{keyres}{Classifying stationary points ($f'(a)=0$)}
\begin{cols}
\textbf{Sign (gradient) test:} check $f'$ just either side of $a$.
\begin{itemize}
  \item $+\ \to\ 0\ \to\ -$: local \textbf{maximum}
  \item $-\ \to\ 0\ \to\ +$: local \textbf{minimum}
  \item $+\to0\to+$ or $-\to0\to-$: \textbf{stationary point of inflection}
\end{itemize}
\nextcol
\textbf{Second derivative test:}
\begin{itemize}
  \item $f''(a)<0$: local maximum (concave down)
  \item $f''(a)>0$: local minimum (concave up)
  \item $f''(a)=0$: \textbf{inconclusive}; use the sign test (e.g.\ $x^4$ has $f''(0)=0$ but a minimum at $0$)
\end{itemize}
\end{cols}
\textbf{Concavity:} $f''(x)>0$ concave up (gradient increasing); $f''(x)<0$ concave down. A \textbf{point of inflection} is where the concavity changes ($f''$ changes sign). There the gradient is a local max or min, so it is the \emph{steepest} point (``the rate of change is greatest/least''). It is a \emph{stationary} point of inflection if also $f'=0$.
\end{keyres}

\begin{defn}{Strictly increasing / decreasing}
$f$ is \textbf{strictly increasing} on an interval $I$ if $x_1<x_2\Rightarrow f(x_1)<f(x_2)$ for all $x_1,x_2\in I$ (strictly decreasing: $f(x_1)>f(x_2)$).
\begin{itemize}
  \item $f'(x)>0$ on an interval $\Rightarrow$ strictly increasing there. The converse fails: $x^3$ is strictly increasing on $\R$ although $f'(0)=0$.
  \item Using the definition, endpoints where $f'=0$ may be included: give the largest intervals, e.g.\ $f(x)=x^2$ is strictly increasing on $[0,\infty)$.
  \item Do \textbf{not} join intervals across an asymptote: $f(x)=\frac1x$ is strictly decreasing on $(-\infty,0)$ and on $(0,\infty)$, but not on $\R\setminus\{0\}$ (since $f(-1)<f(1)$).
\end{itemize}
\end{defn}

\begin{example}{full analysis of a cubic \byhand}
$f(x)=x^3-3x^2-9x+5$. \ $f'(x)=3x^2-6x-9=3(x-3)(x+1)$, \ $f''(x)=6x-6$.\par
\begin{tabular}{@{}l|ccccc@{}}
$x$ & $x<-1$ & $-1$ & $-1<x<3$ & $3$ & $x>3$\\ \hline
sign of $f'(x)$ & $+$ & $0$ & $-$ & $0$ & $+$\\
shape of $f$ & $\nearrow$ & max & $\searrow$ & min & $\nearrow$
\end{tabular}\par
Local maximum $(-1,10)$ [$f''(-1)=-12<0$]; local minimum $(3,-22)$ [$f''(3)=12>0$].\par
Point of inflection: $f''(x)=0\Rightarrow x=1$, and $f''$ changes sign, so $(1,-6)$, with gradient $f'(1)=-12$ (a non-stationary point of inflection).\par
Strictly increasing on $(-\infty,-1]$ and $[3,\infty)$; strictly decreasing on $[-1,3]$.
\end{example}

\begin{example}{a quartic with a stationary point of inflection}
$f(x)=x^4-4x^3$: \ $f'(x)=4x^2(x-3)$. At $x=0$: $f'$ is negative on both sides, so $(0,0)$ is a \textbf{stationary point of inflection}. At $x=3$: local minimum $(3,-27)$.\par
$f''(x)=12x(x-2)=0$ at $x=0$ and $x=2$. Both are points of inflection (sign changes): $(0,0)$ is stationary and $(2,-16)$ is non-stationary.
\end{example}

% -----------------------------------------------------------------
\subsection{Curve sketching with calculus}
\label{sec:sketch}

\begin{cols}[0.53]
\begin{method}{Checklist}
\begin{enumerate}
  \item Domain (maximal or given), endpoints.
  \item Axis intercepts.
  \item $f'(x)=0$ $\to$ stationary points, and their nature.
  \item $f''(x)=0$ with sign change $\to$ points of inflection (if asked).
  \item Asymptotes and behaviour as $x\to\pm\infty$ or near excluded points.
  \item Plot and label everything with coordinates. Use exact values unless told otherwise.
\end{enumerate}
\end{method}
\begin{example}{$f(x)=xe^{-x}$}
Intercept $(0,0)$. \ $f'(x)=e^{-x}-xe^{-x}=e^{-x}(1-x)$, so there is a maximum at $\left(1,\frac1e\right)$.\par
$f''(x)=e^{-x}(x-2)$, so there is a point of inflection at $\left(2,\frac{2}{e^2}\right)$.\par
As $x\to\infty$, $f(x)\to0^+$ (asymptote $y=0$); as $x\to-\infty$, $f(x)\to-\infty$. Range $\left(-\infty,\frac1e\right]$.
\end{example}
\nextcol[0.445]
\centering
\begin{tikzpicture}
\begin{axis}[mygraph, width=6.6cm, height=5.6cm, xmin=-1.0, xmax=6.3, ymin=-1.6, ymax=0.75, xtick=\empty, ytick=\empty]
  \addplot[plotA, domain=-0.65:6.1] {x*exp(-x)};
  \node[pt] at (axis cs:1,0.3679) {}; \node[lab, above] at (axis cs:1,0.3679) {$\left(1,\frac1e\right)$ max};
  \node[pt] at (axis cs:2,0.2707) {}; \node[lab, above right] at (axis cs:2,0.2707) {$\left(2,\frac{2}{e^2}\right)$ inflection};
  \node[pt] at (axis cs:0,0) {};
  \node[lab] at (axis cs:5.4,-0.12) {$y=0$};
\end{axis}
\end{tikzpicture}
\end{cols}

% -----------------------------------------------------------------
\subsection{Absolute maximum and minimum; optimisation}
\label{sec:opt}

\begin{keyres}{Absolute (global) extrema on an interval}
On a closed interval $[a,b]$, the absolute maximum and minimum occur either at stationary points inside the interval \textbf{or at the endpoints}. Evaluate $f$ at all of them and compare. On an open interval, or with asymptotes, an absolute max/min may not exist. A \emph{local} max is not necessarily the absolute max.
\end{keyres}

\begin{example}{endpoint extrema \byhand}
$f(x)=x^3-3x$ on $[-2,3]$. Stationary points: $f'(x)=3x^2-3=0\Rightarrow x=\pm1$, giving $f(-1)=2$, $f(1)=-2$. Endpoints: $f(-2)=-2$, $f(3)=18$.\par
Absolute maximum $18$ (at the endpoint $x=3$); absolute minimum $-2$ (at $x=-2$ and $x=1$). Range $[-2,18]$.
\end{example}

\begin{method}{Optimisation problems}
\begin{enumerate}
  \item Draw a diagram and define the variables (with units).
  \item Write the quantity to optimise. Use the constraint to reduce it to \textbf{one variable}.
  \item State the \textbf{feasible domain} from the context (lengths $>0$, etc.).
  \item Differentiate and solve $=0$. Check the nature (sign test or $f''$) \textbf{and compare with the endpoints}.
  \item Answer the question asked, in context, with units.
\end{enumerate}
\end{method}

\begin{cols}
\begin{example}{open box \byhand}
Squares of side $x$ cm are cut from the corners of a $12\times12$ cm sheet, which is folded into an open box.\par
$V=x(12-2x)^2$, \ $0<x<6$.\par
$V'(x)=(12-2x)^2-4x(12-2x)=(12-2x)(12-6x)$.\par
$V'=0\Rightarrow x=2$ ($x=6$ is outside the domain). $V'$ changes from $+$ to $-$, so this is a maximum: $V(2)=2(8)^2=128$ cm$^3$.
\end{example}
\nextcol
\begin{example}{distance to a curve \examtwo}
Find the minimum distance from the origin to $y=e^x$.\par
$D(x)=\sqrt{x^2+e^{2x}}$. Minimise $D^2$ instead (same $x$): $\frac{d}{dx}(x^2+e^{2x})=2x+2e^{2x}=0$.\par
This can't be solved exactly by hand; CAS gives $x\approx-0.4263$, so $D_{\min}\approx0.7798$.\par
\textbf{Trick:} minimising $D^2$ avoids the square root.
\end{example}
\end{cols}

% -----------------------------------------------------------------
\subsection{Rates of change in context}
\label{sec:rates}

\begin{cols}
\begin{keyres}{Interpreting derivatives}
\begin{itemize}
  \item $\frac{dV}{dt}$ is the rate of change of $V$ with respect to $t$, in units of ($V$ units) per ($t$ unit), e.g.\ L/min.
  \item The sign says increasing or decreasing; the size says how fast.
  \item ``When is it increasing fastest?'' means maximise the rate: solve $\frac{d^2V}{dt^2}=0$ (a point of inflection of $V$).
  \item Chain rule for linked rates: $\frac{dy}{dt}=\frac{dy}{dx}\cdot\frac{dx}{dt}$.
  \item Motion: velocity $v=\frac{dx}{dt}$, acceleration $a=\frac{dv}{dt}$.
\end{itemize}
\end{keyres}
\nextcol
\begin{example}{rate of filling}
$V(t)=1000\left(1-e^{-0.1t}\right)$ litres after $t$ minutes.\par
$V'(t)=100e^{-0.1t}$ L/min, so the initial rate is $V'(0)=100$ L/min.\par
When is the rate $50$ L/min? $e^{-0.1t}=\frac12\Rightarrow t=10\loge2\approx6.93$ min.\par
$V'(t)>0$ always, and $V\to1000$: the tank approaches 1000 L but never reaches it.
\end{example}
\end{cols}

% -----------------------------------------------------------------
\subsection{Newton's method \newsd}
\label{sec:newton}

\begin{cols}[0.53]
\begin{keyres}{The algorithm (on the formula sheet)}
\[x_{n+1}=x_n-\frac{f(x_n)}{f'(x_n)}\]
Start with an estimate $x_0$ of a solution of $f(x)=0$. Each new estimate $x_{n+1}$ is the \textbf{$x$-intercept of the tangent} to $y=f(x)$ at $x=x_n$. Repeat until successive estimates agree to the required accuracy (e.g.\ $|x_{n+1}-x_n|<0.0001$).
\end{keyres}
\begin{warn}[When Newton's method fails or misbehaves]
\begin{itemize}
  \item $f'(x_n)=0$: the tangent is horizontal and never meets the axis (division by zero). E.g.\ $f(x)=x^3-3x+1$ with $x_0=1$.
  \item $f'(x_n)$ near $0$: the next estimate jumps far away, possibly converging to a different root.
  \item The estimates can oscillate (cycle) or diverge. A poor $x_0$ (near a turning point) is the usual cause. Choose $x_0$ close to the root, e.g.\ where $f$ changes sign.
\end{itemize}
\end{warn}
\nextcol[0.445]
\centering
\begin{tikzpicture}
\begin{axis}[mygraph, width=6.4cm, height=5.6cm, xmin=-0.3, xmax=3.4, ymin=-2.6, ymax=8, xtick={1.4621,1.8333,3}, xticklabels={$x_2$,$x_1$,$x_0$}, ytick=\empty]
  \addplot[plotA, domain=0:3.25] {x^2-2};
  \addplot[plotB, thin, domain=1.55:3.15] {7+6*(x-3)};
  \addplot[plotC, thin, domain=1.25:2.15] {1.3611+3.6667*(x-1.8333)};
  \draw[dashed, black!45] (axis cs:3,0) -- (axis cs:3,7);
  \draw[dashed, black!45] (axis cs:1.8333,0) -- (axis cs:1.8333,1.3611);
  \node[pt] at (axis cs:3,7) {};
  \node[pt] at (axis cs:1.8333,1.3611) {};
  \node[pt] at (axis cs:1.4142,0) {};
  \node[lab, above left] at (axis cs:1.40,0.05) {root};
  \node[lab, plotA] at (axis cs:0.55,-1.25) {$y=x^2-2$};
\end{axis}
\end{tikzpicture}

{\scriptsize Each tangent's $x$-intercept is the next estimate: $3\to1.8333\to1.4621\to\dots\to\sqrt2$.}
\end{cols}

\begin{example}{Newton's method}
\textbf{(a) \byhand} Use $x_0=2$ to estimate $\sqrt5$ (a root of $f(x)=x^2-5$, $f'(x)=2x$):\par
$x_1=2-\frac{4-5}{4}=\frac94=2.25$; \quad $x_2=\frac94-\frac{\frac{81}{16}-5}{\frac92}=\frac94-\frac{1}{72}=\frac{161}{72}\approx2.23611$ (compare $\sqrt5\approx2.23607$).\par
\textbf{(b) \examtwo} Solve $x^3-2x-5=0$ with $x_0=2$, where $f'(x)=3x^2-2$:\par
\begin{tabular}{@{}cccc@{}}
$n$ & $x_n$ & $f(x_n)$ & $f'(x_n)$\\ \hline
0 & 2 & $-1$ & 10\\
1 & 2.1 & 0.061 & 11.23\\
2 & 2.094568 & 0.00018 & 11.16\\
3 & 2.094551 & & \\
\end{tabular}\quad The root is $2.0946$ (4 d.p.); successive estimates now agree.
\end{example}
